\subsection{Estratègia d’avaluació i traçabilitat dels requisits}

L’avaluació combina proves funcionals sobre transaccions reals, verificacions qualitatives de les explicacions generades i comprovacions estructurals de l’arquitectura local. No es pretén realitzar una comparació estadística amb productes comercials, sinó demostrar la viabilitat del prototip dins d’un entorn de proves controlat.

La Taula~\ref{tab:tracabilitat_requisits} relaciona els requisits definits al Capítol~\ref{cap:design} amb els escenaris o verificacions utilitzats per comprovar-los. La correspondència és bidireccional: cada escenari de l'apartat~\ref{sec:escenaris_resultats} indica els requisits que valida i, inversament, aquesta taula recull tots els escenaris que intervenen en la validació de cada requisit.

\begin{table}[htbp]
\centering
\small
\renewcommand{\arraystretch}{1.45}
\setlength{\tabcolsep}{9pt}
\begin{tabularx}{\textwidth}{p{0.15\textwidth} | X | p{0.23\textwidth}}
\hline
\textbf{Requisit} & \textbf{Mètode de validació} & \textbf{Escenari o evidència} \\
\hline
RF1 -- Interceptació de transaccions &
Comprovació que el Snap s’activa abans de la confirmació dins de MetaMask. &
Escenaris A i G. \\
\hline
RF2 -- Anàlisi de risc basat en regles &
Execució d’operacions ERC20 i ERC721 amb nivells de risc diferents. &
Escenaris B, C, D, E i G. \\
\hline
RF3 -- Generació d’explicacions comprensibles &
Revisió qualitativa de la coherència, la claredat i la utilitat de l’explicació generada. &
Escenaris B, C, D, E i F. \\
\hline
RF4 -- Integració dins la interfície de la cartera &
Comprovació que l’\textit{insight} es mostra dins de MetaMask abans de confirmar. &
Escenaris A, B, C, D i G. \\
\hline
RNF1 -- Rendiment &
Mesura del temps complet de resposta des de la petició fins a la visualització del resultat. &
Escenari H i, parcialment, G. \\
\hline
RNF2 -- Privadesa &
Verificació que l’anàlisi, la persistència i la generació de text s’executen localment sempre que és possible. &
Escenaris E i G, i revisió arquitectònica. \\
\hline
RNF3 -- Fiabilitat &
Comprovació del comportament del sistema quan el backend o el servei d’IA no estan disponibles. &
Escenari H i, parcialment, A. \\
\hline
RNF4 -- Extensibilitat &
Revisió de la separació modular entre normalització, descodificació, regles, persistència i IA. &
Escenaris E i F, i revisió de la implementació (Capítol~\ref{cap:implementation}). \\
\hline
RNF5 -- Usabilitat &
Valoració qualitativa del llenguatge i de la utilitat de la informació mostrada a l’usuari. &
Escenari F. \\
\hline
\end{tabularx}
\caption{Traçabilitat entre requisits i mecanismes de validació}
\label{tab:tracabilitat_requisits}
\end{table}
